\section{The numerical count gives no obstruction}\label{sec:numerology}

An earlier version of part of this material asserted a numerical obstruction
to secant constructions, on the ground that the Euler characteristic which the
dimension count demands lies below the value at which a linear system on an
abelian variety becomes free of base points.%
\footnote{The content of this section was asserted in the opposite direction in the earlier version cited on the title page, where the two thresholds were presented as necessary conditions and used to rule out a range of constructions. They are sufficient and not necessary, and the two counterexamples below are elementary. We have kept the section rather than deleting it, because an error that was believed is worth printing next to its correction.} That assertion is wrong, and
since it is the kind of error that is easy to repeat we set it out and correct
it here. The corrected conclusion is that the numerical count obstructs
nothing; the obstructions in this paper are the two of
\Cref{thm:divisor,thm:character}, and neither uses any count.

\subsection{The count}

\begin{setupenv}[The secant numerology]\label{setup:numerology}
Let $A$ be of $(K,-1,n)$-Weil type, so $\dim A=\dim\Av=2n$. Suppose $\Av$ is
embedded in $\PP^{N-1}$ by a complete linear system $|M|$ with $M$ ample, so
that $N=h^{0}(M)=\chi(M)$. The $n$-th secant variety of a nondegenerate
$2n$-fold has expected dimension
\[
  \dim\Sec^{n}(\Av)=n(2n+1)-1=2n^{2}+n-1,
\]
and a construction that intersects $\Sec^{n}(\Av)$ with $\Av$ in the expected
dimension inside $\PP^{N-1}$ obtains
\[
  \dim\bigl(\Sec^{n}(\Av)\cap\Av\bigr)
  =\bigl(2n^{2}+n-1\bigr)+2n-\bigl(N-1\bigr)
  =2n^{2}+3n-N .
\]
Asking for this to equal $n$, which is what produces a class of the required
codimension, forces
\begin{equation}\label{eq:chirequired}
  N=\chi(M)=2n^{2}+2n .
\end{equation}
\end{setupenv}

\begin{lemma}\label{lem:chicompare}
For every $n\ge2$,
\[
  2n^{2}+2n<2^{2n}<3^{2n},
\]
and the ratio $2^{2n}/(2n^{2}+2n)$ tends to infinity.
\end{lemma}

\begin{proof}
At $n=2$ the left inequality reads $12<16$. For $n\ge2$ the function
$4^{n}/(2n^{2}+2n)$ has ratio of consecutive terms
$4(2n^{2}+2n)/(2(n+1)^{2}+2(n+1))=4n/(n+2)>1$ for $n>2/3$, so the quotient is
increasing from a value exceeding $1$ at $n=2$; both claims follow.
\end{proof}

\subsection{Why this is not an obstruction}

The theorem of Lefschetz says that for $M$ ample on an abelian variety,
$M^{\otimes2}$ is free of base points and $M^{\otimes3}$ is very ample. If
$M=P^{\otimes2}$ with $P$ ample on a variety of dimension $m$, then
$\chi(M)=2^{m}\chi(P)\ge2^{m}$, and similarly $\chi\ge3^{m}$ in the very ample
case. It is tempting to read these as thresholds: no base-point-free system
below $2^{m}$, no very ample system below $3^{m}$. Both readings invert a
sufficient condition into a necessary one, and both are false already for
abelian surfaces.

\begin{proposition}[The thresholds are not necessary]\label{prop:thresholds}
Let $S$ be a general complex abelian surface, so $m=2$.
\begin{enumerate}[label=\textup{(\roman*)},nosep]
\item A polarisation of type $(1,3)$ on $S$ has $\chi=3<4=2^{2}$ and its
  complete linear system is free of base points; the associated morphism
  $S\to\PP^{2}$ is a six-sheeted covering.
\item A polarisation of type $(1,5)$ on $S$ has $\chi=5<9=3^{2}$ and is very
  ample; the image is a Horrocks--Mumford abelian surface of degree $10$ in
  $\PP^{4}$ \textup{\cite{HM73}}.
\end{enumerate}
In particular neither $2^{m}$ nor $3^{m}$ is a lower bound for the Euler
characteristic of a base-point-free, respectively very ample, ample line
bundle on an abelian variety of dimension $m$.
\end{proposition}

\begin{proof}
Both statements are in the classification of $(1,d)$-polarised abelian
surfaces; see Gross and Popescu \cite[\S1.2]{GP98}, where the base locus of
$|L|$ is determined for each $d$: it is empty for $d\ge3$, consists of four
points for $d=2$, and $|L|$ is very ample for $d\ge5$ on a general surface,
the case $d=5$ being the Horrocks--Mumford embedding of \cite{HM73}. The Euler
characteristics are $\chi=d$ by Riemann--Roch on an abelian surface.
\end{proof}

\begin{corollary}\label{cor:nonumerical}
\Cref{setup:numerology} imposes no contradiction. The requirement
$\chi(M)=2n^{2}+2n$ of \eqref{eq:chirequired} is compatible with $M$ being
very ample for all $n$ for which a polarisation of that Euler characteristic
and suitable type exists on a $2n$-dimensional abelian variety; at $n=2$ the
requirement $\chi=12$ on an abelian fourfold is met by polarisations of type
$(1,1,2,6)$ and $(1,1,1,12)$, and nothing in the numerology excludes very
ampleness for either.
\end{corollary}

\begin{remark}[What we retract]\label{rem:retract}
The claim withdrawn here is the assertion that the polarisation required by
\eqref{eq:chirequired} admits no base-point-free linear system. It appeared in
the supplementary material of the earlier version cited on the title page,
together with a table of values of $2n^{2}+2n$, $2^{2n}$ and $3^{2n}$. The
table is arithmetically correct and \Cref{lem:chicompare} reproduces it; the
inference drawn from it is not. We are grateful to have found this before it
propagated, and we note that the obstructions this paper does prove,
\Cref{thm:divisor,thm:character}, are of a different character: they bound
where a class can lie in a decomposition of cohomology, and no choice of
polarisation, ambient space or object affects them.
\end{remark}
